% Part: lambda-calculus
% Chapter: introduction
% Section: composition

\documentclass[../../../include/open-logic-section]{subfiles}

\begin{document}

\olfileid{lam}{int}{com}
\olsection{\usetoken{S}{lambda definable} অপেক্ষকগুলি মিশ্রণের অধীনে বদ্ধ}

\begin{lem}
!!{lambda definable} অপেক্ষকগুলি মিশ্রণের অধীনে বদ্ধ।
\end{lem}

\begin{proof}
ধরা যাক $h$, $g_0,$ \dots,~$g_{k-1}$ থেকে মিশ্রণের সাহায্যে~$f$-কে
সংজ্ঞায়িত করা হয়েছে। যথাক্রমে $H$, $G_0$, \dots,~$G_{k-1}$ যদি
$h$, $g_0$, \dots,~$g_{k-1}$-কে !!{lambda defined} করে, তাহলে আমাদের
এমন একটি পদ~$F$ খুঁজতে হবে, যা~$f$-কে !!{lambda define}s। কিন্তু আমরা
স্রেফ নিচের মতো করে~$F$-কে সংজ্ঞায়িত করতে পারি:
\[
F(x_0, \dots, x_{l-1}) = H(G_0(x_0, \dots, x_{l-1}),
\dots, G_{k-1}(x_0, \dots, x_{l-1})).
\]
অন্যভাবে বললে, ল্যাম্বডা ক্যালকুলাসের ভাষা মিশ্রণ প্রকাশের পক্ষে
বিশেষভাবে উপযোগী।
\end{proof}

\end{document}

